\section{A SAMPLE OF TEXT WITH TABLE}\label{table}
%This sample section presents a table and allows LaTeX to create a
 %list of tables for the template.  The original table appeared in
 %Eduardo's thesis.  The surrounding text does not provide enough
 %explanation to understand what the table presents.
%The next line produces an indented paragraph to start the document
 %unit.  The LaTeX defaults start most units without indentations.
\hspace{\parindent}
Now we want to describe a tricomplex $B{\mathcal D}^3$
so that its total differential corresponds to the total differential $d''$ in the
version of $\Tot ({\mathcal D}^3)$ modified by deleting all the
vertical  $b'$-complexes from the tricomplex ${\mathcal D}^3$.
Let us consider the tricomplex $B{\mathcal D}^3$ as a tridimensional array
that results after replacing all the vertical $b'$-complexes by zero complexes.
The rules determining the maps connecting these modules are described in the following table.
\begin{table}[h]
\caption{Table of Maps}
%Table 1. Table of Maps
\begin{center}
\begin{tabular}{|c|c|c|c|} \hline
Case  
               & To determine                     
                                     & Look at              
                                             &   In       
 \\ \hline \hline
(1) 
               & $(B{\mathcal D}^3)_{p,q,0} \rightarrow (B{\mathcal D}^3)_{p+q-j-1,j,0}$  
                                    & $(j{+}1,q{+}1)\text{-entry of}$   
                                             &  $d'_{p+q}$          
\\ 
               & $0\leq j\leq p+q-1$ 
                                    & $B_{p+q,p+q+1}$               
                                             &                    
\\ \hline
(2) 
     & $(B{\mathcal D}^3)_{p,q,0}\rightarrow  (B{\mathcal D}^3)_{p+q-j-2,j,1}$  
               & $(j{+}1,q{+}1)\text{-entry of}$  
                                             &  $d'_{p+q}$          
\\ 
    & 
               & $B_{p+q-1,p+q+1}$             
                                             &                    
\\ \hline
(3) 
     & $(B{\mathcal D}^3)_{p,q,1}\rightarrow  (B{\mathcal D}^3)_{p+q-j,j,0}$    
               & $(j{+}1,q{+}1)\text{-entry of}$   
                                              &  $d'_{p+q+1}$        
\\ 
    & 
    
               & $B_{p+q+1,p+q+1}$             &                    
\\ \hline
(4) 
     & $(B{\mathcal D}^3)_{p,q,1}\rightarrow  (B{\mathcal D}^3)_{p+q-j-1,j,1}$  
             & $(j{+}2,q{+}1)\text{-entry of}$   
                                               &  $d'_{p+q+1}$        
\\
    &  
             & $B_{p+q,p+q+1}$                 &                    
\\ \hline
\end{tabular}
\end{center}
\end{table}

Let us start with case (1) in Table 1.
%This is the end of the file table.tex
